\documentclass[11pt,a4paper]{article}
\usepackage[margin=2.5cm]{geometry}
\usepackage[colorlinks=true,linkcolor=blue,urlcolor=blue]{hyperref}
\usepackage{booktabs}

\title{The Agent Stack:\\Scripts, Programs, Prompts, Skills, Plugins, Agents}
\author{Omar}
\date{September 2026}

\begin{document}
\maketitle

\section{The model I started with}

\begin{verbatim}
1. Script      (deterministic)
2. Program     (application, bigger than script)
3. Skill       (contains scripts and instructions)
4. Plugin      (contains scripts, instructions, and other things)
5. Agent       (the thing that controls all these)
\end{verbatim}

Two things in it are right: deterministic code sits below instruction
packages, and the agent is the thing in control. One thing is wrong: this is
a single stack. It is not. The \emph{artifacts} stack; the agent does not
belong on the stack, it stands \emph{beside} it. And containment runs in
every direction, so the whole system is a graph, not a tower.

\section{The fix: two axes and one runtime}

Every item is one of three kinds:

\begin{enumerate}
  \item \textbf{Deterministic code} (script, program): given input X,
    produces output Y. No judgment. If it works once, it works forever.
  \item \textbf{Instruction packages} (prompt, skill, plugin): text plus
    optional resources that tell a model what to do. Inert until an agent
    reads them.
  \item \textbf{Runtime} (agent): a model in a loop, with tools. The only
    participant that decides anything.
\end{enumerate}

\section{Corrected order, low to high}

\begin{table}[h]
\centering
\small
\begin{tabular}{@{}c l l p{6.2cm}@{}}
\toprule
\textbf{\#} & \textbf{Thing} & \textbf{Kind} & \textbf{Notes} \\
\midrule
1 & Script & code & interpreted glue, no build step \\
2 & Program & code & compiled binary; a real software distinction, but
  invisible to an agent: both levels 1--2 are just ``callable deterministic
  code'' \\
3 & Prompt & instruction (raw) & \textbf{new level, inserted here.} Needs a
  model to interpret, unlike code. No packaging, no trigger, one-shot: dies
  with the conversation \\
4 & Skill & instruction (packaged) & directory: SKILL.md plus optional
  scripts/, references/, assets/. A prompt plus resources plus a trigger;
  reusable; progressively disclosed \\
5 & Plugin & instruction (bundled) & directory: skills/, commands/, agents/,
  hooks/, MCP config. A distribution format: several skills plus harness
  extensions \\
-- & Agent & runtime & agent .md definition plus a running process. Not
  above the others as a container: beside them as the reader and executor \\
\bottomrule
\end{tabular}
\caption{Physical forms: script/program are files or binaries; prompt is
text; skill and plugin are directories; an agent is a definition file while
it is waiting, and a process while it runs.}
\end{table}

Why the prompt sits between program and skill:

\begin{itemize}
  \item A script or program needs \textbf{no model}: that is why they are
    below the line.
  \item A prompt \textbf{needs a model but no packaging}: one-shot text.
  \item A skill needs a model \textbf{and} comes with packaging: trigger,
    resources, reusability.
\end{itemize}

\section{Containment: the stack is really a graph}

The second realization was correct: anything can contain anything.

\begin{itemize}
  \item \textbf{program contains agent}: any program that calls an LLM API
    inside a loop. Claude Code itself is a level-2 program (a compiled CLI)
    hosting a level-5 agent. The stack loops on itself.
  \item \textbf{skill contains script contains agent}: a skill's script may
    itself call an LLM API; now a skill contains an agent.
  \item \textbf{skill points to skill}: a skill body can tell the agent to
    invoke another skill.
  \item \textbf{plugin contains skills, agents, commands, hooks}.
  \item \textbf{CLAUDE.md} is a standing prompt; a \textbf{slash command} is
    a saved prompt with a manual trigger.
\end{itemize}

So each level is a containment \emph{option}, not a law. From any single
node the view is a tree; the whole system is a directed graph. What puts the
agent ``on top'' is not containment but \textbf{control}: the agent is the
only thing in the picture that decides.

\section{Do plugins contain agents?}

Yes. A plugin can carry an \texttt{agents/} directory:

\begin{verbatim}
my-plugin/
|-- agents/
|   `-- reviewer.md      <- subagent definition
|-- skills/
|-- commands/
`-- hooks/
\end{verbatim}

\texttt{reviewer.md} is a markdown file with YAML frontmatter (name,
description, tools, model) and a body of instructions. The harness
registers it; the main agent can spawn it as a subagent: a fresh model
instance with a fresh context, running that definition. The recursion is
real too: a plugin loaded by the main agent's harness can define subagents
that the main agent then spawns. The plugin stays inert the whole time; the
subagent is a live process.

\section{Do agents exist physically, like SKILL.md?}

In two senses:

\begin{itemize}
  \item \textbf{Definition}: yes. Literally a \texttt{.md} file with
    frontmatter (\texttt{\textasciitilde/.claude/agents/foo.md}, or a
    plugin's \texttt{agents/foo.md}). As inert as SKILL.md.
  \item \textbf{Instance}: no. A running process: a model holding that
    definition as its instructions, plus tools, plus a task. Exists only
    while the loop runs. You can \texttt{cat} the definition; you can never
    \texttt{cat} the instance.
\end{itemize}

The symmetry:

\begin{verbatim}
SKILL.md  : instructions FOR an existing agent
agent.md  : instructions TO HATCH a new agent
\end{verbatim}

Both inert until the harness uses them. The harness treats them
differently: one is loaded into an existing context, the other spawns a new
one.

\section{Where the prompt fits in the flow}

\begin{verbatim}
You type a prompt (level 3)
  -> harness assembles context: system prompt + CLAUDE.md + skill
     descriptions
  -> agent reads it, decides
  -> agent runs a script or program (levels 1-2)
  -> or loads a skill body (level 4), possibly from an installed
     plugin (level 5)
  -> or spawns a subagent (another agent .md)
  -> observes the result, repeats
\end{verbatim}

The prompt is the \textbf{input event} that starts every cycle. Everything
else on the instruction axis is a way of storing prompts so they survive:
CLAUDE.md (always-on prompt), slash command (manual prompt), skill
(self-triggering prompt plus resources), agent definition (a prompt that
becomes a whole new mind).

\section{Summary}

\begin{itemize}
  \item \textbf{Code} (script, program): deterministic, decides nothing.
  \item \textbf{Instructions} (prompt, skill, plugin): text plus resources,
    decides nothing: packaging at increasing levels of structure.
  \item \textbf{Agent}: the only decider: a model in a loop.
  \item \textbf{Containment} is arbitrary in every direction;
    \textbf{control} is not: everything inert is waiting for an agent to
    read or run it.
\end{itemize}

\end{document}